\documentclass[12pt]{article}
\usepackage[utf8]{inputenc}
\usepackage[T1]{fontenc}
\usepackage{hyperref}
\title{Fania Language Resources for K-12 Curriculum Development}
\author{Compiled by Claude Sonnet 5.5 (claude-sonnet-5-5) for Andr\'as Kornai}
\date{}
\begin{document}
\maketitle

\section{Introduction}
Fania (self-name Kulaale; also Fanya, Fagnia, Kobe; ISO 639-3 \texttt{fni}) is an endangered Adamawa (Bua group) language of the southern Gu\'era region of Chad, with roughly 1,000 speakers. Search results show only academic and reference material; no educational material was found.

\section{Primary Educational Resources}
None found. No primers, textbooks or school programmes appeared in the search results.

\section{Linguistic and Reference Materials}
\begin{itemize}
\item Lionnet, ``Kulaale (Chad)'' Language Documentation \& Description snapshot (2020): \url{https://www.lddjournal.org/article/id/1213/}
\item Same snapshot on HAL: \url{https://halshs.archives-ouvertes.fr/halshs-04370663}
\item Glottolog record for the Lionnet reference: \url{https://glottolog.org/resource/reference/id/580904}
\item Numerals of Fania (Max Planck channumerals database): \url{https://lingweb.eva.mpg.de/channumerals/Fania.htm}
\item Wikipedia: \url{https://en.wikipedia.org/wiki/Fania_language}
\end{itemize}

\section{Religious and Community Resources}
Joshua Project lists the Fanya as a Muslim, unreached group; a Bible translation is not reported: \url{https://joshuaproject.net/people_groups/11189}

\section{Digital and Audio Resources}
None found beyond the documentation snapshot above.

\section{Curriculum Development Context}
The language is spoken in about a dozen villages, and young people who move to towns such as Sarh and N'Djamena often do not acquire it. No orthography or literacy material was found, so any curriculum would start from documentation and community work, building on the Lionnet description.
\end{document}